% Options for packages loaded elsewhere
\PassOptionsToPackage{unicode}{hyperref}
\PassOptionsToPackage{hyphens}{url}
\PassOptionsToPackage{dvipsnames,svgnames,x11names}{xcolor}
\documentclass[
  10pt,
  fontset=fandol]{ctexart}
\usepackage{xcolor}
\usepackage[margin=16mm]{geometry}
\usepackage{amsmath,amssymb}
\setcounter{secnumdepth}{-\maxdimen} % remove section numbering
\usepackage{iftex}
\ifPDFTeX
  \usepackage[T1]{fontenc}
  \usepackage[utf8]{inputenc}
  \usepackage{textcomp} % provide euro and other symbols
\else % if luatex or xetex
  \usepackage{unicode-math} % this also loads fontspec
  \defaultfontfeatures{Scale=MatchLowercase}
  \defaultfontfeatures[\rmfamily]{Ligatures=TeX,Scale=1}
\fi
\usepackage{lmodern}
\ifPDFTeX\else
  % xetex/luatex font selection
\fi
% Use upquote if available, for straight quotes in verbatim environments
\IfFileExists{upquote.sty}{\usepackage{upquote}}{}
\IfFileExists{microtype.sty}{% use microtype if available
  \usepackage[]{microtype}
  \UseMicrotypeSet[protrusion]{basicmath} % disable protrusion for tt fonts
}{}
\makeatletter
\@ifundefined{KOMAClassName}{% if non-KOMA class
  \IfFileExists{parskip.sty}{%
    \usepackage{parskip}
  }{% else
    \setlength{\parindent}{0pt}
    \setlength{\parskip}{6pt plus 2pt minus 1pt}}
}{% if KOMA class
  \KOMAoptions{parskip=half}}
\makeatother
\setlength{\emergencystretch}{3em} % prevent overfull lines
\providecommand{\tightlist}{%
  \setlength{\itemsep}{0pt}\setlength{\parskip}{0pt}}
\usepackage{amsmath,amssymb,mathtools,needspace}
\xeCJKDeclareCharClass{CJK}{"2032,"0400->"04FF,"2160->"216B,"2460->"2473,"25A0,"25C7,"25CB,"25CF}
\IfFontExistsTF{Arial Unicode MS}{\xeCJKsetup{AutoFallBack=true}\setCJKfallbackfamilyfont{\CJKrmdefault}{Arial Unicode MS}}{}
\setcounter{secnumdepth}{0}
\setlength{\emergencystretch}{2em}
\allowdisplaybreaks
\usepackage{bookmark}
\IfFileExists{xurl.sty}{\usepackage{xurl}}{} % add URL line breaks if available
\urlstyle{same}
\hypersetup{
  pdftitle={偶阶求积公式的代数精度},
  colorlinks=true,
  linkcolor={Maroon},
  filecolor={Maroon},
  citecolor={Blue},
  urlcolor={Blue},
  pdfcreator={LaTeX via pandoc}}

\title{偶阶求积公式的代数精度}
\author{}
\date{}

\begin{document}
\maketitle

\subsubsection{4.2.2 偶阶求积公式的代数精度}\label{ux5076ux9636ux6c42ux79efux516cux5f0fux7684ux4ee3ux6570ux7cbeux5ea6}

作为插值型的求积公式，\(n\) 阶的 Newton-Cotes 公式至少具有 \(n\) 次的代数精度（见定理 4.1）。实际的代数精度能否进一步提高呢？

先看 Simpson 公式（4.2.3），它是二阶 Newton-Cotes 公式，因此至少具有二次代数精度。进一步用 \(f(x)=x^3\) 进行检验，按 Simpson 公式计算，得

\[
S=\frac{b-a}{6}\left[a^3+4\left(\frac{a+b}{2}\right)^3+b^3\right].
\]

另一方面直接求积得 \(I=\int_a^b x^3\,\mathrm{d}x=\dfrac{b^4-a^4}{4}\)。这时有 \(S=I\)，即 Simpson 公式对次数不超过三次的多项式均能准确成立。又容易验证它对于 \(f(x)=x^4\) 通常是不准确的，因此，Simpson 公式实际上具有三次代数精度。

一般地，可以证明下述论断。

\textbf{定理 4.2} 当阶 \(n\) 为偶数时，Newton-Cotes 公式（4.2.1）至少有 \(n+1\) 次代数精度。

\textbf{证明} 只要验证，当 \(n\) 为偶数时，Newton-Cotes 公式对 \(f(x)=x^{n+1}\) 的余项为零。

按余项公式（4.1.7），由于这里 \(f^{(n+1)}(x)=(n+1)!\)，从而有

\[
R[f]=\int_a^b\prod_{j=0}^n(x-x_j)\,\mathrm{d}x.
\]

引进变换 \(x=a+th\)，并注意到 \(x_j=a+jh\)，有

\[
R[f]=h^{n+2}\int_0^n\prod_{j=0}^n(t-j)\,\mathrm{d}t.
\]

若 \(n\) 为偶数，则 \(\dfrac n2\) 为整数，再令 \(t=u+\dfrac n2\)，进一步有

\[
R[f]=h^{n+2}\int_{-n/2}^{n/2}\prod_{j=0}^n\left(u+\frac n2-j\right)\,\mathrm{d}u,
\]

\href{../../source-images/pdf-098.jpeg}{原页图像}

据此可以断定 \(R[f]=0\)，因为被积函数

\[
H(u)=\prod_{j=0}^n\left(u+\frac n2-j\right)=\prod_{j=-n/2}^{n/2}(u-j)
\]

是个奇函数。证毕。

\begin{center}\rule{0.5\linewidth}{0.5pt}\end{center}

\subsection{相关原页校注（agent 补充）}\label{ux76f8ux5173ux539fux9875ux6821ux6ce8agent-ux8865ux5145}

以下校注来自本节覆盖的原页，可能也涉及同页相邻小节；原文照录与校正说明分开保留。

\subsubsection{原书 PDF 页 97 的校注}\label{ux539fux4e66-pdf-ux9875-97-ux7684ux6821ux6ce8}

\href{../pages/pdf-097.md}{完整原页转录} · \href{../../source-images/pdf-097.jpeg}{原页图像}

校注记录：CH04A-E01。

\begin{quote}
校注（agent 补充，CH04A-E01）：上述``当 \(n\geq8\) 时''若理解为对所有这类整数 \(n\) 成立，则存在原书表述错误。保留原文。由式（4.2.2）直接积分，\(n=9\) 时系数为 \[
\left(\frac{2857}{89600},\frac{15741}{89600},\frac{27}{2240},\frac{1209}{5600},\frac{2889}{44800},\frac{2889}{44800},\frac{1209}{5600},\frac{27}{2240},\frac{15741}{89600},\frac{2857}{89600}\right),
\] 全部为正；\(n=8\) 出现负系数这一事实仍成立。上述反例为 agent 按原公式进行的符号积分核验，不属于原书正文。
\end{quote}

\subsection{本节覆盖原页的校注记录}\label{ux672cux8282ux8986ux76d6ux539fux9875ux7684ux6821ux6ce8ux8bb0ux5f55}

下列记录按原页关联，含同页相邻小节的校注；计算时同时读取上文校注和完整原页。

\begin{itemize}
\tightlist
\item
  \texttt{CH04A-E01}：原书 PDF 页 97
\end{itemize}

\end{document}
